\secciongrande{Mecanismos de ejecución de código en un navegador web la arquitectura del navegador}

\resumen{
  La interacción en un navegador inicia con la solicitud de un recurso y finaliza con su visualización, adaptándose a propósitos como velocidad, fidelidad o seguridad. La arquitectura de referencia, compartida por la mayoría de navegadores, se compone de módulos clave como la interfaz de usuario, el motor del buscador, el motor de renderizado, el subsistema de comunicaciones, el intérprete de JavaScript, el parser XML, el componente de visualización y el subsistema de persistencia de datos.
}

\definicion{Arquitectura de referencia}{Conjunto de componentes básicos comunes a la mayoría de los navegadores web, que definen su funcionamiento esencial.}

\subtituloA{Proceso de interacción del usuario}

\lista{
  \textbf{Inicio y fin:} La interacción comienza al indicar una dirección de recurso y finaliza con la visualización del contenido. |
  \textbf{Dependencias:} El proceso se configura según el propósito (velocidad, fidelidad al contenido, seguridad de la comunicación).
}